% SHARD-ID: AQM-49-REGULAR-FUNCTIONS-AFFINE-LINE-WEYL-LABELS
% SHARD-TITLE: Regular Functions, Supports, and Weyl Labels on the Affine Line
% SHARD-SUMMARY: Separates regular functions as residue profiles from finite-support Weyl labels.
% SHARD-SUMMARY: Defines polynomial zero sets as finite closed supports and proves the Chinese-remainder independence of their qudits.
% SHARD-SUMMARY: Distinguishes reduced smearing over \(Z(h)\) from multiplicity-sensitive Artin-ring thickening.
% SHARD-KEYWORDS: regular function, affine line, closed support, Chinese remainder, Weyl label, thickening

\section{Regular Functions, Supports, and Weyl Labels on the Affine Line}
\label{sec:regular-functions-affine-line-weyl-labels}

\subsection{Status and Source Anchors}
This shard fixes the current meaning of ``regular function'' in the
closed-point Weyl-net layer.  It refines AQM-28 and AQM-29.  The source-local
facts are Stacks: affine-scheme global sections in
\path{references/algebraic_geometry/stacks_project_schemes.tex}, lines
670--719; \(\Spec\), \(D(f)\), and \(V(I)\) in
\path{references/algebraic_geometry/stacks_project_algebra.tex}, lines
2972--3023 and 3104--3211; residue fields in the same file, lines 250--274
and 3310--3350.  Finite-field Weyl operators are AQM-28, sourced from
Ashikhmin--Knill as recorded there.  The Artin-ring layer is AQM-36.

\subsection{Lamport Dependency Skeleton}
\[
\begin{array}{rcl}
D1&:&k=\mathbb F_Q,\quad R=k[x],\quad X=\Spec R,\quad
      \mathcal P=\{\pi\in R:\pi\text{ monic irreducible}\}.\\
D2&:&x_\pi=(\pi),\quad K_\pi=R/(\pi).\\
D3&:&h\ne0,\quad D(h)=\{\mathfrak p:h\notin\mathfrak p\},\quad
      \Gamma(D(h),\mathcal O_X)=R_h.\\
D4&:&\operatorname{ev}_U(F,G)_\pi=(F\bmod\pi,G\bmod\pi)
      \in K_\pi^2,\quad x_\pi\in U.\\
L1&:&\text{A nonzero regular-function pair on nonempty }D(h)
      \text{ has cofinite closed support in }D(h).\\
T1&:&\operatorname{ev}_U(F,G)\in E(U)\text{ iff }(F,G)=(0,0).\\
D5&:&Z_{\rm cl}(h)=\{x_\pi:\pi\mid h\},\quad
      \operatorname{rad}(h)=\prod_{\pi\mid h}\pi.\\
L2&:&Z_{\rm cl}(h)\text{ is finite and depends only on }
      \operatorname{rad}(h).\\
D6&:&E_h^{\rm red}=\bigoplus_{\pi\mid h}K_\pi^2,\quad
      W_h^{\rm red}(F,G)=W_{Z_{\rm cl}(h)}((F\bmod\pi,G\bmod\pi)_\pi).\\
L3&:&W_h^{\rm red}(F,G)\text{ depends only on }F,G\bmod\operatorname{rad}(h).\\
T2&:&\text{Every finite closed-support Weyl label is represented by
      polynomial pairs modulo the squarefree support polynomial.}\\
D7&:&h=\prod_i\pi_i^{e_i}\text{ also defines the thickened rings }
      R/(\pi_i^{e_i}).\\
T3&:&\text{Multiplicities }e_i\text{ are invisible in }E_h^{\rm red}
      \text{ and visible only in the Artin-ring layer.}
\end{array}
\]

\subsection{Regular Functions as Residue Profiles D1--T1}
Let \(k=\mathbb F_Q\), \(R=k[x]\), and \(X=\Spec R\).  By the Stacks affine
source,
\[
  \Gamma(X,\mathcal O_X)=R.
\]
More generally, on the nonempty standard open \(D(h)\), \(h\ne0\), the regular
functions are the localized ring
\[
  \Gamma(D(h),\mathcal O_X)=R_h.
\]
For \(h=1\), this is the global case.

Let \(U=D(h)\ne\emptyset\).  A closed point of \(U\) is \(x_\pi\) with
\(\pi\nmid h\).  For \(F=u/h^m\in R_h\), define
\[
  F(x_\pi)=(u\bmod\pi)(h\bmod\pi)^{-m}\in K_\pi .
\]
This is well-defined because \(h\bmod\pi\ne0\) in the field \(K_\pi\).
For a pair \((F,G)\in R_h^2\), define the residue profile
\[
  \operatorname{ev}_U(F,G)_\pi=(F(x_\pi),G(x_\pi))\in K_\pi^2.
\]

\paragraph{Lemma \(L1\).}
If \((F,G)\ne(0,0)\) in \(R_h^2\), then the set of closed points
\[
  \{\pi:\pi\nmid h,\ \operatorname{ev}_U(F,G)_\pi\ne(0,0)\}
\]
is cofinite in \(\{\pi:\pi\nmid h\}\).

\emph{Proof.}
Write \(F=u/h^m\) and \(G=v/h^n\).  The profile is zero at \(x_\pi\in U\)
exactly when \(u\bmod\pi=0\) and \(v\bmod\pi=0\).  This means that \(\pi\)
divides both \(u\) and \(v\).  If \((F,G)\ne(0,0)\), then at least one of
\(u,v\) is nonzero after localization.  A nonzero polynomial has only finitely
many monic irreducible divisors.  Hence only finitely many \(\pi\nmid h\) can
make both residues zero.

\paragraph{Theorem \(T1\).}
Under the current closed-point Weyl-net convention,
\[
  \operatorname{ev}_U(F,G)\in E(U)
  \quad\Longleftrightarrow\quad
  (F,G)=(0,0).
\]

\emph{Proof.}
By AQM-28,
\[
  E(U)=\bigoplus_{x_\pi\in U}K_\pi^2
\]
means finite closed support.  If \((F,G)=(0,0)\), the profile has empty support.
If \((F,G)\ne(0,0)\), Lemma \(L1\) gives cofinite support in the infinite
closed-point set of \(U\); AQM-29 proves that \(D(h)\) omits only finitely many
closed points from the infinite set \(\mathcal P\).  Hence the profile is not
finite-support and is not an element of \(E(U)\).

Thus a regular function is a classical global or local residue profile.  It is
not, by itself, a quasi-local Weyl observable.

\subsection{Polynomials as Finite Closed Supports D5--L2}
Now let \(h\in R\) be nonzero and factor
\[
  h=c\prod_{i=1}^r\pi_i^{e_i},
  \qquad c\in k^\times,\quad e_i\ge1,
\]
with distinct monic irreducibles \(\pi_i\).  Define
\[
  Z_{\rm cl}(h)=\{x_{\pi_i}:1\le i\le r\},
  \qquad
  \operatorname{rad}(h)=\prod_{i=1}^r\pi_i .
\]

\paragraph{Lemma \(L2\).}
\(Z_{\rm cl}(h)\) is finite and depends only on \(\operatorname{rad}(h)\), not
on the exponents \(e_i\).

\emph{Proof.}
A closed point \(x_\pi\) lies in \(V(h)\) exactly when \(h\in(\pi)\), which is
equivalent to \(\pi\mid h\).  The factorization of a nonzero polynomial has
only finitely many irreducible factors.  Replacing \(h\) by
\(\operatorname{rad}(h)\) preserves the same divisibility relation
\(\pi\mid h\).

This is the first quantum use of a polynomial: it selects a finite closed
support.  It does not yet choose a Weyl operator.

\subsection{Reduced Smearing over a Support D6--T2}
Define the reduced support label space
\[
  E_h^{\rm red}=\bigoplus_{\pi\mid h}K_\pi^2.
\]
For polynomial pairs \((F,G)\in R^2\), define the reduced smeared label
\[
  e_h^{\rm red}(F,G)=((F\bmod\pi,\ G\bmod\pi))_{\pi\mid h}\in E_h^{\rm red}
\]
and the reduced smeared Weyl operator
\[
  W_h^{\rm red}(F,G)=W_{Z_{\rm cl}(h)}(e_h^{\rm red}(F,G)).
\]

\paragraph{Lemma \(L3\).}
\(e_h^{\rm red}(F,G)\), hence \(W_h^{\rm red}(F,G)\), depends only on the
classes of \(F\) and \(G\) modulo \(\operatorname{rad}(h)\).

\emph{Proof.}
The label at \(x_\pi\) is obtained by reducing modulo \(\pi\).  Two polynomials
have the same reduction modulo every \(\pi\mid h\) iff their difference is
divisible by every \(\pi\mid h\).  Since these irreducibles are pairwise
coprime, this is equivalent to divisibility by \(\operatorname{rad}(h)\).

\paragraph{Theorem \(T2\).}
For every finite closed support \(S=\{x_{\pi_1},\ldots,x_{\pi_r}\}\) and every
label
\[
  e=((q_i,p_i))_{i=1}^r\in\bigoplus_{i=1}^rK_{\pi_i}^2,
\]
there exist polynomials \(F,G\in R\) with
\[
  F\bmod\pi_i=q_i,\qquad G\bmod\pi_i=p_i
  \quad\text{for every }i.
\]
Thus distinct closed-point qudits have no finite compatibility condition in
the reduced Weyl net.

\emph{Proof.}
The ideals \((\pi_i)\) are pairwise coprime.  The Chinese remainder theorem
therefore gives a surjection
\[
  R\longrightarrow\prod_{i=1}^r R/(\pi_i)=\prod_{i=1}^rK_{\pi_i}.
\]
Choose \(F\) mapping to \((q_i)_i\) and \(G\) mapping to \((p_i)_i\).  These
polynomials have the required residues.  Since every finite label is obtained
this way, the qudit factors indexed by distinct irreducibles are independent
inside the finite-support Weyl algebra.

\subsection{Multiplicity and the Artin-Ring Layer D7--T3}
The same polynomial \(h=\prod_i\pi_i^{e_i}\) may also be read as a closed
subscheme with multiplicity.  Then the local rings
\[
  A_i=R/(\pi_i^{e_i})
\]
are finite local Artin rings supported at the same closed points \(x_{\pi_i}\).
AQM-36 quantizes this thickened layer by
\[
  E_h^{\rm thick}=\bigoplus_i A_i^2,\qquad
  \mathcal H_h^{\rm thick}=\bigotimes_i\ell^2(A_i).
\]

\paragraph{Theorem \(T3\).}
The reduced support operator \(W_h^{\rm red}(F,G)\) forgets all multiplicities
\(e_i\).  Multiplicities become quantum degrees of freedom only after replacing
the residue fields \(K_{\pi_i}\) by the Artin rings \(A_i\).

\emph{Proof.}
By Lemma \(L3\), the reduced label only uses \(F,G\) modulo
\(\operatorname{rad}(h)=\prod_i\pi_i\).  Therefore \(h\) and
\(\operatorname{rad}(h)\) give the same reduced label space, the same Hilbert
space \(\bigotimes_i\ell^2(K_{\pi_i})\), and the same reduced Weyl operators.
AQM-36 instead uses \(R/(\pi_i^{e_i})\), whose cardinality is
\(Q^{e_i\deg\pi_i}\).  These larger rings record nilpotent jet coefficients at
the same closed point.  They are not additional Zariski points.

\subsection{Conclusion}
In the present framing, regular functions have three separate roles:
\[
\begin{array}{rcl}
\text{profile} &:& (F,G)\in\Gamma(U,\mathcal O_X)^2
        \mapsto \operatorname{ev}_U(F,G),\text{ usually infinite support};\\
\text{support selector} &:& h\ne0\mapsto Z_{\rm cl}(h),\text{ a finite closed
        support};\\
\text{smeared operator} &:& (h,F,G)\mapsto W_h^{\rm red}(F,G),
        \text{ or }W_h^{\rm thick}(F,G)\text{ after Artin thickening}.
\end{array}
\]
Only the last line is a Weyl observable, and only after a finite support or a
thickened finite scheme has been chosen.
